% 2.1.2
% Ficheiro incluído por main.tex

\subsection[Sincronização do relógio (chrony)]{Sincronização do relógio (chrony)}

\esquema{2-1-2-chrony}{Sincronização do relógio por NTP}

\textbf{Objetivo:} acertar a hora do Ubuntu pela do DC. O Kerberos rejeita pedidos com mais de \textbf{5 minutos} de diferença entre o cliente e o servidor.

\passo{Passo 1 - Direcionar o chrony ao Windows Server}

\begin{lstlisting}[style=bashstyle]
grep -q "^server <IP_DC>" /etc/chrony/chrony.conf || \
  echo "server <IP_DC> iburst" | sudo tee -a /etc/chrony/chrony.conf
grep "^server <IP_DC>" /etc/chrony/chrony.conf
\end{lstlisting}

Só acrescenta se ainda não existir, por isso não duplica. \texttt{server} = fonte de tempo; \texttt{iburst} = sincroniza mais depressa no arranque.

\passo{Passo 2 - Reiniciar o chrony e confirmar}

\begin{lstlisting}[style=bashstyle]
sudo systemctl restart chrony
systemctl status chrony --no-pager
chronyc sources -v
\end{lstlisting}

\emph{Verificação esperada:} \texttt{active (running)} (se o serviço não existir, tenta \texttt{chronyd}) e, no \texttt{chronyc sources -v}, a linha de \texttt{\textless{}IP\_DC\textgreater{}} com \texttt{\textasciicircum{}*} ou \texttt{\textasciicircum{}+} - prova de que o relógio está mesmo a ser acertado pelo DC. Se \texttt{\textless{}IP\_DC\textgreater{}} aparecer com \texttt{\textasciicircum{}?}, espera uns segundos e repete o comando.

\problemas{\hyperref[ec-b1-kerberos]{Erros Comuns - 2.1 Relógio e Kerberos}}
